% \section{}

\vspace{-0.2cm}
\section{Limitations, future work and broader impact}
\vspace{-0.2cm}

Our framework offers a novel approach to address differentially private training, but also introduces new challenges. \rebuttal{Clipless DP-SGD replaces a bias of optimizer with a bias in function space (as discussed in Appendix~\ref{sec:discussionbias}). We primary rely on GNP networks, where high performing architectures are quite different from the usual CNN architectures. This serves as \textbf{(1) a motivation to further develop Gradient Norm Preserving architectures}.} As emphasized in Remark~\ref{remark:gnp}, we anticipate that progress in these areas will greatly enhance the effectiveness of our approach.  
Additionally, to meet requirement~\ref{req:pgd}, we rely on projections, necessitating additional efforts to incorporate ``trivializations''~\cite{trockman2021orthogonalizing,singla2021skew}. %recent advancements associated with differentiable reparametrizations~\cite{trockman2021orthogonalizing,singla2021skew}.
Finally, as mentioned in Remark~\ref{remark:forwardbounds}, our propagation bound method can be refined. \rebuttal{Comparison of Clipless DP-SGD with recent methods based on back-propagation, in particular~\cite{bu2023differentially}, is discussed in appendix~\ref{sec:comparisonsota}}. %Beyond the application of Algorithms~\ref{alg:gnbc}~and~\ref{alg:noclipDP-SGDlocal}, our framework provides numerous opportunities to enhance our understanding of prevalent techniques identified in the literature. % An in-depth exploration of these is beyond the scope of this work, so we focus on giving insights on promising tracks based on our theoretical analysis.
Furthermore, the development of networks with known Lipschitz constant with respect to parameters is a question of independent interest, making \lstinline[basicstyle=\ttfamily, backgroundcolor=\color{gray!20}]{lip-dp} \textbf{(2) a useful tool for the study of the optimization dynamics} in Lipschitz constrained neural networks.


